\documentclass[12pt]{article}
\usepackage[utf8]{inputenc}
\usepackage[T1]{fontenc}
\usepackage{hyperref}
\title{Yakama Language Resources for K-12 Curriculum Development}
\author{Compiled by Claude Sonnet 5.5 (claude-sonnet-5-5) for Andr\'as Kornai}
\date{}
\begin{document}
\maketitle

\section{Introduction}
Yakama (Yakima Sahaptin, Ichish\-k\'iin) is a Sahaptian language of Washington State. A recent news feature says only a few dozen fluent speakers remain. The search found a published dictionary and a university teaching-materials archive, but no ready-made K-12 curriculum.

\section{Primary Educational Resources}
\begin{itemize}
\item UW Sahaptin Teaching Materials Archive, which preserves Virginia Beavert's lesson materials, including a Yakama beginner's guide, handouts and class recordings. Some pages are password-protected: \url{https://depts.washington.edu/sahteach/tag/yakama/page/3/}
\item Same archive, vocabulary category: \url{https://depts.washington.edu/sahteach/category/vocabulary/page/20/}
\end{itemize}

\section{Linguistic and Reference Materials}
\begin{itemize}
\item \emph{Ichishk\'iin S\textbackslash i\'nwit Yakama: Yakima Sahaptin Dictionary} (Beavert and Hargus, University of Washington Press and Heritage University, 2010), as described in the UBC RelLex entry: \url{https://wiki.ubc.ca/Documentation:RelLex/Ichishk%C3%ADin_S%C9%A8%CC%81nwit_Yakama:_Yakima_Sahaptin_Dictionary}
\end{itemize}

\section{Religious and Community Resources}
None were found in the search.

\section{Digital and Audio Resources}
\begin{itemize}
\item Class recordings in the UW Sahaptin archive (tag: recording): \url{https://depts.washington.edu/sahteach/tag/recording/}
\end{itemize}

\section{Curriculum Development Context}
\begin{itemize}
\item Stanford Magazine feature on a learner studying the language in Toppenish, which describes the very small number of fluent speakers: \url{https://stanfordmag.org/contents/found-in-translation}
\end{itemize}
The search found no Yakama Nation school programme or standardized curriculum, so they would need to be checked directly. Because few speakers remain, recordings and the dictionary are the key assets for any curriculum.
\end{document}